\scanpage{225}
\mspage{- 58 -}
n\textsuperscript{o} 215

sorte que $J^T_{y,x}$ est l'application réciproque de $J^T_{x,y}$) ; par ailleurs, $\xi$ étant déterminé comme ci-dessus, on a $s_{x,\xi}(z)=J^T_{y,z}(\handwritten{\eta})$ pour tout $z\in U$, de sorte que $\handwritten{\eta}\unsure{\handwritten{\longrightarrow}|\longrightarrow} s_{x,\xi}(z)$ est un isomorphisme de $T_y^T$ sur $T_{\unsure{z|x}}^T$ ; les sections $s_{x,\xi}$ forment donc bien un débitage. De plus, il résulte immédiatement que ce débitage ne dépend pas du choix de $x$ dans $U$. \unsure{Nobs|Nous} l'appellerons le \uline{débitage de} $T^T$ \uline{associé à} $D$, et nous le désignerons par $D^T$. Ceci dit, on définit une structure de fût sur $T^T$ en prenant comme ensemble fondamental de débitages l'ensemble des débitages associés à ceux de $E$.

Soit maintenant $T'$ une espèce de structures associée à $T$ et telle que, si $H$ est une structure d'espèce $T$ sur un ensemble $\ne\handwritten{\varnothing}$, l'ensemble de base de la structure d'espèce $T'$ associée à $H$ soit non vide. On peut alors considérer $T'$ comme une espèce de structures associée à $S$ (par l'intermédiaire \unsure{\uline{de}|de} $T$). Soit $J$ un isomorphisme d'une structure $F$ d'espèce $S$ sur une structure $F'$ de même espèce ; soient $H$ et $H'$ les structures d'espèce $T$ associées à $F$, $F'$ et $K$, $K'$ les structures d'espèce $T'$ associées à $H$, $H'$ (donc aussi à $F$, $F'$, quand on considère $T'$ comme associée à $S$). L'isomorphisme $J$ définit un isomorphisme $J^T$ de $H$ sur $H'$, et ce dernier un isomorphisme $(J^T)^{T'}$ de $K$ sur $K'$. Il est clair que $(J^T)^{T'}$ est identique à l'isomorphisme $J^{T'}$ de $K$ sur $K'$ qui se déduit de $J$ quand on considère $T'$ comme associée à $S$. Si donc $E$ est un fût d'espèce $S$, $E^T$ le fût associé d'espèce $T$ et $(E^T)^{T'}$ le fût d'espèce $T'$ associé à $E^T$, $(E^T)^{T'}$ est le fût d'espèce $T'$ associé à $E$ ; c'est la transitivité des fûts associés.

PROPOSITION 2. - \uline{Soient} $E$ \uline{un fût d'espèce} $S$, $T$ \uline{une espèce de}
% archive sheet number: 224
